% !TeX program = xelatex
\documentclass[12pt,a4paper]{article}

\usepackage{fontspec}
\usepackage{polyglossia}
\setmainlanguage{russian}
\setotherlanguage{english}
\setmainfont{FreeSerif}[
  BoldFont       = FreeSerifBold,
  ItalicFont     = FreeSerifItalic,
  BoldItalicFont = FreeSerifBoldItalic
]
\setmonofont{DejaVu Sans Mono}[Scale=0.82]

\usepackage{longtable}
\usepackage{threeparttable}   % gives ThreePartTable + TableNotes, for use with longtable
\usepackage{pdflscape}        % landscape page for the wide table

\usepackage[left=3cm,right=1.5cm,top=2cm,bottom=2cm]{geometry}
\usepackage{setspace}
\usepackage{indentfirst}
\usepackage{titlesec}
\usepackage{listings}
\usepackage{xcolor}
\usepackage{colortbl}
\usepackage{enumitem}
\usepackage{booktabs}
\usepackage{amsmath}
\usepackage{graphicx}
\usepackage{array}
\usepackage{amssymb}
\usepackage{microtype}
\usepackage[hidelinks]{hyperref}

% Spacing
\setlength{\parindent}{1.25cm}
\setlength{\parskip}{4pt plus 1pt}
\onehalfspacing

% Section headings
\titleformat{\section}
  {\normalfont\fontsize{14}{17}\bfseries}
  {\thesection.}{0.6em}{}
\titlespacing*{\section}{0pt}{20pt plus 3pt}{10pt plus 2pt}

\titleformat{\subsection}
  {\normalfont\fontsize{12}{15}\bfseries}
  {\thesubsection}{0.6em}{}
\titlespacing*{\subsection}{0pt}{14pt plus 2pt}{6pt plus 1pt}

% Colours
\definecolor{codebg}{RGB}{250,250,250}
\definecolor{codeframe}{RGB}{210,210,210}
\definecolor{kwcolor}{RGB}{0,0,160}
\definecolor{commentcolor}{RGB}{60,128,60}
\definecolor{strcolor}{RGB}{160,30,30}
\definecolor{tblred}{RGB}{160,0,0}

% Code listings
\lstdefinelanguage{plpgsql}{
  morekeywords={CREATE,OR,REPLACE,FUNCTION,RETURNS,AS,BEGIN,END,
    DECLARE,SELECT,INTO,FROM,WHERE,IF,THEN,ELSE,ELSIF,RETURN,
    LANGUAGE,plpgsql,RAISE,EXCEPTION,NOT,FOUND,IS,NULL,AND,UPDATE,
    SET,INSERT,DELETE,OLD,NEW,TG_OP,INTEGER,FLOAT,VARCHAR,DATE,
    TRIGGER,AFTER,ON,FOR,EACH,ROW,EXECUTE,ALTER,TABLE,ADD,COLUMN,
    MATERIALIZED,VIEW,REFRESH,COUNT,DEFAULT,DISTINCT},
  sensitive=false,
  morecomment=[l]{--},
  morestring=[b]',
}

\lstset{
  language          = plpgsql,
  basicstyle        = \small\ttfamily\setstretch{1.1},
  keywordstyle      = \color{kwcolor}\bfseries,
  commentstyle      = \color{commentcolor}\itshape,
  stringstyle       = \color{strcolor},
  backgroundcolor   = \color{codebg},
  frame             = single,
  framesep          = 4pt,
  rulecolor         = \color{codeframe},
  breaklines        = true,
  breakatwhitespace = true,
  columns           = flexible,
  postbreak         = \mbox{\textcolor{codeframe}{$\hookrightarrow$}\space},
  showstringspaces  = false,
  tabsize           = 4,
  xleftmargin       = 8pt,
  xrightmargin      = 8pt,
  aboveskip         = 10pt,
  belowskip         = 8pt,
}

\newcolumntype{L}[1]{>{\raggedright\arraybackslash}p{#1}}

% Allow slightly loose lines to avoid overfull hboxes from long \texttt tokens
\setlength{\emergencystretch}{4em}
\hfuzz=3pt
\hbadness=10000

\begin{document}

% ════════════════════════════════════════════════════════
%  TITLE PAGE
% ════════════════════════════════════════════════════════
% ========== TITLE PAGE ==========
\begin{titlepage}
    \centering
    \textbf{Санкт-Петербургский Национальный Исследовательский Университет ИТМО}\\
    \textbf{Факультет Программной Инженерии и Компьютерной Техники}

    \vspace{2cm}
    \includegraphics[width=0.4\textwidth]{ITMOLogo.png}\\[1cm]

    \vspace{3cm}
    Лабораторная работа №6\\[0.4em]
    По {\normalsize <<Основы профессиональной деятельности>>}\\[16pt]
    {\large Вариант №~738}

    \vfill
    \begin{flushright}
        Выполнил студент:\\
        Эллити Мохамед Эмад Ахмед Авад\\[1em]
        Группы:\\ P3131 \\[1em]
        Преподаватель:\\
Клименков Сергей Викторович \\
Зайцева Ирина Сергеевна
    \end{flushright}

    \vspace{2cm}
    Санкт-Петербург 2026 г.
\end{titlepage}

\tableofcontents
\clearpage

% ════════════════════════════════════════════════════════
%  ТЕКСТ ЗАДАНИЯ
% ════════════════════════════════════════════════════════
\section{Текст задания}

По выданному преподавателем варианту разработать и исследовать работу комплекса программ обмена данными в режиме прерывания программы. Основная программа должна изменять содержимое заданной ячейки памяти (X), которое должно быть представлено как знаковое число. Область допустимых значений изменения X должна быть ограничена заданной функцией F(X) и конструктивными особенностями регистра данных ВУ (8-ми битное знаковое представление). Программа обработки прерывания должна выводить на ВУ модифицированное значение X в соответствии с вариантом задания, а также игнорировать все необрабатываемые прерывания.

\begin{enumerate}
    \item Основная программа должна увеличивать на 3 содержимое X (ячейки памяти с адресом $020_{16}$) в цикле.
    \item Обработчик прерывания должен по нажатию кнопки готовности ВУ-1 осуществлять вывод результата вычисления функции $F(X) = -3X-8$ на данное ВУ, а по нажатию кнопки готовности ВУ-2 выполнить операцию побитового «ИЛИ» содержимого регистра данных данного ВУ и X, результат записать в X.
    \item Если X оказывается вне ОДЗ при выполнении любой операции по его изменению, то необходимо в X записать минимальное по ОДЗ число.
\end{enumerate}

% ════════════════════════════════════════════════════════
%  ОПИСАНИЕ ПРОГРАММЫ
% ════════════════════════════════════════════════════════
\section{Описание программы}

\subsection*{Назначение программы}

\begin{enumerate}
    \item Основная программа должна увеличивать на 3 содержимое X (ячейки памяти с адресом $020_{16}$) в цикле.
    \item Обработчик прерывания должен по нажатию кнопки готовности ВУ-1 осуществлять вывод результата вычисления функции $F(X) = -3X-8$ на данное ВУ, а по нажатию кнопки готовности ВУ-2 выполнить операцию побитового «ИЛИ» содержимого регистра данных данного ВУ и X, результат записать в X.
    \item Если X оказывается вне ОДЗ при выполнении любой операции по его изменению, то необходимо в X записать минимальное по ОДЗ число.
\end{enumerate}

\subsection*{Текст программы}

\begin{lstlisting}[language={[x86masm]Assembler}]
        ORG     0x000
V0:     WORD    $default, 0x180
V1:     WORD    $int1,    0x180     ; вектор ВУ-1
V2:     WORD    $int2,    0x180     ; вектор ВУ-2
V3:     WORD    $default, 0x180
V4:     WORD    $default, 0x180
V5:     WORD    $default, 0x180
V6:     WORD    $default, 0x180
V7:     WORD    $default, 0x180

        ORG     0x020
X:      WORD    ?
max:    WORD    0x0028              ; 40, максимум ОДЗ
min:    WORD    0xFFD3              ; -45, минимум ОДЗ

default:
        IRET                        ; необрабатываемые прерывания

START:  DI
        CLA
        OUT     0x1                 ; запрет прерываний для неиспользуемых ВУ
        OUT     0x7
        LD      #0x9                ; MR: IE=1, вектор 1
        OUT     3
        LD      #0xA                ; MR: IE=1, вектор 2
        OUT     5
        EI

main:   DI                          ; атомарность на время изменения X
        LD      X
        ADD     #3
        CMP     min
        BGE     chk1_max
        JUMP    chk1_ldmin
chk1_max:
        CMP     max
        BLT     chk1_store
        BEQ     chk1_store
chk1_ldmin:
        LD      min
chk1_store:
        ST      X
        EI
        JUMP    main

int1:   DI                          ; ВУ-1: вывод F(X) = -3X-8
        PUSH
        LD      X
        ADD     X                   ; 2X
        ADD     X                   ; 3X
        NEG                         ; -3X
        SUB     #8                  ; -3X-8
        OUT     2
        POP
        EI
        IRET

int2:   DI                          ; ВУ-2: X := X | РД(ВУ-2)
        PUSH
        CLA
        IN      4
        OR      X
        CMP     min
        BGE     chk2_max
        JUMP    chk2_ldmin
chk2_max:
        CMP     max
        BLT     chk2_store
        BEQ     chk2_store
chk2_ldmin:
        LD      min
chk2_store:
        ST      X
        POP
        EI
        IRET
\end{lstlisting}

% ════════════════════════════════════════════════════════
%  ОДЗ
% ════════════════════════════════════════════════════════
\section*{Область допустимых значений}

\[
-128 \le -3X-8 \le 127
\]
\[
-120 \le -3X \le 135
\]
\[
-45 \le X \le 40
\]

\[
40 = 0000.0000.0010.1000 = 0x0028
\]
\[
-45 = 1111.1111.1101.0011 = 0xFFD3
\]

Число $X \in [\text{FFD3};\ 0028]$.

% ════════════════════════════════════════════════════════
%  РАСПОЛОЖЕНИЕ ДАННЫХ
% ════════════════════════════════════════════════════════
\section*{Расположение данных в памяти}

\begin{itemize}
    \item Вектор прерываний: 0x000 -- 0x00F
    \item Переменные (X, max, min): 0x020 -- 0x022
    \item Программа: 0x023 – 0x054
\end{itemize}

% ════════════════════════════════════════════════════════
%  ОБЛАСТЬ ПРЕДСТАВЛЕНИЯ
% ════════════════════════════════════════════════════════
\section*{Область представления}

\begin{itemize}
    \item X, min, max -- знаковое 16-разрядное целое число;
    \item DR КВУ -- 8-ми разрядное целое знаковое число.
\end{itemize}

% ════════════════════════════════════════════════════════
%  ВЫВОД
% ════════════════════════════════════════════════════════
\section*{Вывод}

В ходе выполнения лабораторной работы я изучил обмен данными с ВУ-1 и ВУ-2 в режиме прерываний, разобрался с циклом прерывания и порядком сохранения/восстановления состояния процессора при входе и выходе из обработчика. Также закрепил знания в написании программ на ассемблере БЭВМ.

% ════════════════════════════════════════════════════════
%  МЕТОДИКА ПРОВЕРКИ
% ════════════════════════════════════════════════════════
\section*{Методика проверки программы}

\textbf{Проверка обработки прерываний:}

\begin{enumerate}
    \item Загрузить текст программы в БЭВМ.
    \item Заменить \texttt{NOP} по нужному адресу на \texttt{HLT} (либо использовать точку останова).
    \item Запустить программу в режиме РАБОТА.
    \item Установить «Готовность ВУ-1».
    \item Дождаться останова.
    \item Записать текущее значение X из памяти БЭВМ (адрес 0x020).
    \item Записать результат обработки прерывания -- содержимое DR контроллера ВУ-1 (порт \#2).
    \item Рассчитать ожидаемое значение $F(X)=-3X-8$ и сравнить с результатом.
    \item Нажать «Продолжение».
    \item Ввести в ВУ-2 произвольное число, установить «Готовность ВУ-2».
    \item Дождаться останова, записать значение X, сравнить с ожидаемым $X\,|\,\text{РД(ВУ-2)}$ (с учётом коррекции по ОДЗ).
    \item Нажать «Продолжение».
\end{enumerate}

\textbf{Проверка основной программы:}

\begin{enumerate}
    \item Загрузить текст программы в БЭВМ.
    \item Записать в переменную X значение, близкое к максимуму ОДЗ (например, 37).
    \item Запустить программу в пошаговом режиме.
    \item Пройти нужное количество шагов, убедиться, что при превышении X значения 40 происходит сброс в минимальное по ОДЗ значение (-45).
\end{enumerate}

\begin{table}[h!]
\centering
\small
\begin{tabular}{|c|c|c|}
\hline
\multicolumn{3}{|c|}{\textbf{Прерывание ВУ-1}} \\
\hline
\textbf{X (AC)} & \textbf{Ожидание $F(X)$} & \textbf{DR ВУ-1} \\
\hline
$0A_{16}$ (10) & $DA_{16}$ (-38) & $DA_{16}$ (-38) \\
$D3_{16}$ (-45) & $7F_{16}$ (127) & $7F_{16}$ (127) \\
$28_{16}$ (40) & $80_{16}$ (-128) & $80_{16}$ (-128) \\
\hline
\end{tabular}
\end{table}

\begin{table}[h!]
\centering
\small
\begin{tabular}{|c|c|c|}
\hline
\multicolumn{3}{|c|}{\textbf{Прерывание ВУ-2}} \\
\hline
\textbf{X до} & \textbf{DR ВУ-2} & \textbf{X после ($X\,|\,\text{DR}$, с коррекцией ОДЗ)} \\
\hline
$05_{16}$ (5) & $0A_{16}$ (10) & $0F_{16}$ (15) \\
$28_{16}$ (40) & $0F_{16}$ (15) & $D3_{16}$ (-45) -- коррекция (47$>$40) \\
$D3_{16}$ (-45) & $01_{16}$ (1) & $D3_{16}$ (-45) \\
\hline
\end{tabular}
\end{table}

\begin{table}[h!]
\centering
\small
\begin{tabular}{|c|c|}
\hline
\multicolumn{2}{|c|}{\textbf{Основная программа ($X:=X+3$)}} \\
\hline
\textbf{X до} & \textbf{X после} \\
\hline
$25_{16}$ (37) & $28_{16}$ (40) \\
$28_{16}$ (40) & $D3_{16}$ (-45) -- коррекция (43$>$40) \\
$D3_{16}$ (-45) & $D6_{16}$ (-42) \\
\hline
\end{tabular}
\end{table}

\end{document}